\begin{proposition}[Quadratic collapse of the toric Mellin coefficient]
\label{prop:toric-mellin-square}
Let $p\geq5$ be prime, let $\chi$ be the quadratic character of
$\mathbb F_p$ extended by $\chi(0)=0$, and let
\[
 \mu_p(a)=\#\{(x,y,z)\in(\mathbb F_p^\times)^3:
                   \Lambda(x,y,z)=a\}
 \qquad(a\in\mathbb F_p^\times),
\]
where $\Lambda$ is the Laurent polynomial in
\eqref{eq:oracleC-Lambda}.  Put
\[
 A_j(t)=\sum_{k=0}^j\binom jk^2t^k
\]
and, for $1\leq j\leq p-2$, put
\[
 S_{p,j}(t)=\sum_{r\in\mathbb F_p^\times}r^j
 \chi\bigl((t-1)^2r^2-2(t+1)r+1\bigr).
\]
Then the fiber formula below is an integer equality, and the remaining
identities hold in~$\mathbb F_p$.

\begin{enumerate}
\item For every $a\ne0$,
  \begin{equation}
   \mu_p(a)=(p-2)^2+
   \sum_{\substack{t,w\in\mathbb F_p^\times\\
                    w\ne1,\ tw\ne1}}
   \chi\!\left(
    \left(t+1-\frac{t(w-1)(1-tw)}{aw}\right)^2-4t
   \right).
   \label{eq:toric-fiber-quadratic}
  \end{equation}

\item The inner quadratic Mellin transform collapses exactly:
  \begin{equation}
   \boxed{S_{p,j}(t)=-A_{p-1-j}(t).}
   \label{eq:toric-inner-collapse}
  \end{equation}
  Consequently
  \begin{equation}
   \boxed{
   b_j=-\sum_{t\in\mathbb F_p^\times}
       t^jA_j(t)A_{p-1-j}(t).}
   \label{eq:toric-Mellin-pairing}
  \end{equation}

\item If $1\leq j\leq(p-1)/2$, then the polynomial identity
  \begin{equation}
   A_{p-1-j}(t)=(1-t)^{p-1-2j}A_j(t)
   \qquad\text{in }\mathbb F_p[t]
   \label{eq:toric-Legendre-duality}
  \end{equation}
  turns~\eqref{eq:toric-Mellin-pairing} into the weighted-square formula
  \begin{equation}
   \boxed{
   b_j=-\sum_{t\in\mathbb F_p^\times}
       t^j(1-t)^{p-1-2j}A_j(t)^2.}
   \label{eq:toric-weighted-square}
  \end{equation}
\end{enumerate}
\end{proposition}

\begin{proof}
On a nonzero fiber of~$\Lambda$, the change of variables
\[
 u=\frac{1+y}{y},\qquad
 v=\frac{1+z}{z},\qquad
 w=-\frac{x}{uv}
\]
is invertible and has
\[
 u,v,w\in\mathbb F_p^\times,qquad
 u,v\ne1,qquad w\ne1,qquad uvw\ne1.
\]
Writing $t=uv$ and $s=u+v$, direct substitution gives
\[
 \Lambda=\frac{t(w-1)(1-tw)}{w(u-1)(v-1)},
 \qquad
 (u-1)(v-1)=t-s+1.
\]
Thus $\Lambda=a$ determines
\[
 s=t+1-\frac{t(w-1)(1-tw)}{aw}.
\]
For fixed $t,w$, the number of ordered pairs $(u,v)$ with product~$t$
and sum~$s$ is $1+\chi(s^2-4t)$.  The conditions $u,v\ne1$ are automatic:
the excluded equality $s=t+1$ would force $(w-1)(1-tw)=0$.
There are
\[
 (p-2)+(p-2)(p-3)=(p-2)^2
\]
admissible pairs $(t,w)$; the first summand is the slice $t=1$, where the
two excluded values of~$w$ coincide.  This proves
\eqref{eq:toric-fiber-quadratic}.

Set
\[
 Q_t(R)=(t-1)^2R^2-2(t+1)R+1.
\]
The classical Legendre generating identity is
\begin{equation}
 Q_t(R)^{-1/2}=\sum_{n\geq0}A_n(t)R^n.
 \label{eq:toric-Legendre-generating}
\end{equation}
Since $Q_t(0)=1$, fractional powers are unambiguously defined in
$\mathbb F_p[[R]]$.  Frobenius gives
\[
 Q_t(R)^{(p-1)/2}
 =Q_t(R)^{-1/2}\bigl(Q_t(R)^p\bigr)^{1/2}
 =Q_t(R)^{-1/2}Q_t(R^p)^{1/2}
 \equiv Q_t(R)^{-1/2}\pmod{R^p}.
\]
Euler's criterion and multiplicative orthogonality therefore give
\begin{align*}
 S_{p,j}(t)
 &=\sum_{r\ne0}r^jQ_t(r)^{(p-1)/2}\\
 &=-[R^{p-1-j}]Q_t(R)^{(p-1)/2}
  =-A_{p-1-j}(t),
\end{align*}
which proves~\eqref{eq:toric-inner-collapse}.

For $1\leq j\leq p-2$, torus orthogonality and the exponent range in
Lemma~\ref{lem:oracleC-constant-term} give
\begin{equation}
 \sum_{a\ne0}\mu_p(a)a^j
 =\sum_{(x,y,z)\in(\mathbb F_p^\times)^3}\Lambda(x,y,z)^j
 =-b_j.
 \label{eq:toric-fiber-Mellin}
\end{equation}
The constant term $(p-2)^2$ in~\eqref{eq:toric-fiber-quadratic} has zero
$j$-th Mellin coefficient.  For
\[
 \kappa(t,w)=\frac{t(w-1)(1-tw)}w,
\]
multiplication by $a^2$ inside the quadratic character followed by the
bijection $a=\kappa r$ gives
\[
 \sum_{a\ne0}a^j
 \chi\!\left(\left(t+1-\frac\kappa a\right)^2-4t\right)
 =\kappa^jS_{p,j}(t).
\]
Because $j\geq1$, the values $w=1,t^{-1}$ may now be restored: their
$\kappa^j$ is zero.  If
\[
 L_t(w)=\frac{(w-1)(1-tw)}w=1+t-tw-w^{-1},
\]
then every exponent in $L_t(w)^j$ lies in $[-j,j]$, so
\[
 \sum_{w\ne0}\kappa(t,w)^j
 =-t^j\operatorname{CT}_w L_t(w)^j
 =-t^jA_j(t).
\]
Combining this with~\eqref{eq:toric-inner-collapse} and
\eqref{eq:toric-fiber-Mellin} proves
\eqref{eq:toric-Mellin-pairing}.

Finally, if $P_n$ denotes the Legendre polynomial, then
\[
 A_n(t)=(1-t)^nP_n\!\left(\frac{1+t}{1-t}\right).
\]
The hypergeometric expression
$P_n(X)={}_2F_1(-n,n+1;1;(1-X)/2)$ shows, for
$0\leq j\leq(p-1)/2$, that
\[
 P_{p-1-j}(X)=P_j(X)\qquad\text{in }\mathbb F_p[X]:
\]
the two upper parameters are the same unordered pair modulo~$p$, and the
series truncates at degree~$j$.  Clearing the powers of $1-t$ proves the
polynomial identity~\eqref{eq:toric-Legendre-duality}, including $t=1$.
Substitution in~\eqref{eq:toric-Mellin-pairing} gives
\eqref{eq:toric-weighted-square}.
\end{proof}

\begin{remark}[The collapse is an exact reformulation, not yet a saving]
\label{rem:toric-Mellin-tautology}
Expanding~\eqref{eq:toric-Mellin-pairing} and summing over~$t$ gives
\[
 -\sum_{t\ne0}t^jA_j(t)A_{p-1-j}(t)
 =\sum_{\substack{0\leq k\leq j,\ 0\leq\ell\leq p-1-j\\
                   j+k+\ell=p-1}}
   \binom jk^2\binom{p-1-j}{\ell}^2.
\]
Putting $\ell=p-1-j-k$ and using
\[
 \binom{p-1-j}{k}\equiv(-1)^k\binom{j+k}{k}\pmod p
\]
recovers exactly
\[
 b_j=\sum_{k=0}^j\binom jk^2\binom{j+k}{k}^2.
\]
For $k>p-1-j$, both sides vanish modulo~$p$ by the same carry.  Thus the
quadratic fiber sum, the Legendre pairing, and the original Ap\'ery
coefficient are coefficientwise equivalent.  In particular,
\eqref{eq:toric-weighted-square} is a square in an indefinite finite-field
pairing, not a positive norm; it can vanish, for example at $(p,j)=(11,5)$.
Any zero-set bound still requires arithmetic information beyond this exact
collapse.
\end{remark}
